\clearpage\begingroup
\hypertarget{native-271-II}{}
\HistoricalHead{266}{L}{PRAKTISCHE UND SPHÄRISCHE ASTRONOMIE. VERÖFFENTLICHUNGEN UND BRIEFWECHSEL.}
\def\Bsec#1#2{#1\mathord{\overset{\mathrm{s}}{,}}#2}
\def\BCeres{\mathord{\tikz[baseline=-.2ex,x=1em,y=1em,line width=.045em]{
  \draw (.40,.47) .. controls (.30,.66) and (.02,.58) .. (.04,.36) .. controls (.06,.20) and (.25,.18) .. (.35,.29);
  \draw (.32,.34)--(.32,-.02);
  \draw (.19,.10)--(.46,.10);}}}
\def\BJuno{\mathord{\tikz[baseline=-.2ex,x=1em,y=1em,line width=.045em]{
  \draw (.23,.72)--(.23,-.04);
  \draw (.05,.63)--(.41,.39) (.05,.39)--(.41,.63);
  \draw (.06,.13)--(.40,.13);}}}
\begin{center}
[II.]\\[.6em]
[BEOBACHTUNGEN VON PLANETEN.]\\[.7em]
\textsc{Gauss} an \textsc{Olbers}.
\end{center}
\begin{center}\small
\rule{.94\linewidth}{.35pt}\\
{}[\textsc{Wilhelm Olbers}, Sein Leben und seine Werke, II\textsubscript{1}, 1900, S. 201—202, 234, 236, 251.]\\
\rule{.94\linewidth}{.35pt}
\end{center}
\begin{center}Braunschweig, 7. September 1804. — [Beobachtungen der Ceres.]\end{center}
\noindent . . . . . . . . . . . .\par
\noindent\hspace*{\parindent}Endlich habe ich auch zum erstenmale das Vergnügen gehabt, unsre $\BCeres$ hier in Braunschweig zu beobachten. Den 29. August ging sie 32 Ceti um $11^{\mathrm h}18^{\mathrm m}2^{\mathrm s}$ nach einem Mittel aus 3 gut harmonirenden Beobb. $\Bsec{23}5$ vor und war um $11^{\mathrm h}31^{\mathrm m}11^{\mathrm s}$ um $1'3''$ nördlicher. Daraus finde ich
\begin{center}\small\setlength{\tabcolsep}{6pt}
\begin{tabular}{rlrl}
&&&Differenz von d. Ephem.\\
$11^{\mathrm h}18^{\mathrm m}2^{\mathrm s}$ t. m.&AR app.&$15^0\;0'\,32''$&$-7'\,32''$\\
&Decl. app.&$9^0\,55'\,23''$ austr.&$+4'\;3''$.
\end{tabular}
\end{center}
\clearpage\SourceLeaf{272}\HistoricalHead{267}{R}{BEOBACHTUNGEN VON PLANETEN IN BRAUNSCHWEIG.}
\noindent\hspace*{\parindent}. . . . . . Vorgestern habe ich die $\BCeres$ wieder beob[achtet]; allein die Beobb., welche unterbrochen wurden, stimmten schlecht. Nach einer folgte sie
\begin{center}um $11^{\mathrm h}1^{\mathrm m}$ auf 28 Ceti $22^{\mathrm s}$ und ging vor 30 voraus $78^{\mathrm s}$\end{center}
\noindent nach der $2^{\mathrm{ten}}$
\begin{center}um $11^{\mathrm h}16^{\mathrm m}$ auf 28 Ceti $19^{\mathrm s}$ und ging voraus vor $30:81^{\mathrm s}$.\end{center}
\noindent Decl. aus beiden Beobb. im Mittel $10^0\,36'\,52''$. Doch ist diese noch zweifelhafter, da der Declinationsunterschied gerade sehr unvortheilhaft für den Durchmesser des Gesichtsfeldes war. . . . . . .
\begin{center}Braunschweig, 23. November 1804. — [Beobachtungen der Juno.]\end{center}
\noindent . . . . . . . . . . . .\par
\noindent\hspace*{\parindent}Ich selbst habe die Juno noch am 16., 18., 19. beobachtet und mit dem $\ast$ 8. Grösse, der nördl[ich] folgte, verglichen. Die Beobachtungen sind aber alle nicht sehr zuverlässig, besonders die Declination
\begin{center}\small\setlength{\tabcolsep}{4pt}\renewcommand{\arraystretch}{1.15}
\begin{tabular}{r rrr rrr rrr l}
&\multicolumn{3}{c}{M[ittlere] Z[eit]}&&&&&&&\\
Nov. 16.&$6^{\mathrm h}$&$39^{\mathrm m}$&$56^{\mathrm s}$&$356^0$&$17'$&$58''$&$10^0$&$59'$&$25''$&\\
18.&6&21&6&356&33&56&10&58&41&\\
19.&6&4&43&356&42&47&10&58&34&wohl zu gross
\end{tabular}
\end{center}
\noindent u[nd] nur auf Eine Beobachtung gegründet. . . . . . .
\begin{center}Braunschweig, 7. December 1804. — [Beobachtung der Ceres.]\end{center}
\noindent . . . . . . . . . . . .\par
\noindent\hspace*{\parindent}Ceres habe ich am 4. Dec[ember] gleichfalls beobachtet:
\begin{center}$7^{\mathrm h}10^{\mathrm m}\qquad 2^0\,48'\,9''\qquad 10^0\,57'\,45''.$\end{center}
\noindent Die Ephemeride gibt also die AR $6'$ zu klein, die Decl. $3\tfrac34{}'$ zu gross. . . . .
\begin{center}Braunschweig, 7. December 1804. — [Beobachtung der Juno.]\end{center}
\noindent\hspace*{\parindent}Seit meinem letzten Briefe . . . sind mir noch gar keine auswärtige Beobb. der Juno wieder zugekommen; Ihre vom 18. Nov. und die Seeberger vom 6. Nov. waren also die letzten. Da ich indess wünschte, meinen astronomischen Freunden eine verbesserte Ephemeride mittheilen zu können, so habe ich
\HistoricalFoot{}{34\textsuperscript{*}}
\clearpage\SourceLeaf{273}\HistoricalHead{268}{L}{PRAKTISCHE UND SPHÄRISCHE ASTRONOMIE. VERÖFFENTLICHUNGEN UND BRIEFWECHSEL.}
\noindent am 4. Dec. die $\BJuno$ einmal wieder selbst beobachtet und mit 17 Bode und dem darauf folgenden $\ast$ 7. 8. Grösse wiederholt verglichen. Die Position dieser Sterne habe ich selbst aus der \textit{Hist. Cél.} abgeleitet und damit gefunden:
\begin{center}\small
Dec. 4.\quad $6^{\mathrm h}43^{\mathrm m}$ M. Z.\quad AR $359^0\,42'\,17''$\quad D. $10^0\,16'\,3''$.
\end{center}
\noindent Die Beobachtungen stimmten nicht schlecht, doch fällt der etwanige Fehler in der Grösse des Gesichtsfeldes ganz auf die Declination . . . . . .
\begin{center}Braunschweig, 12. Februar 1805. — [Beobachtungen der Juno.]\end{center}
\noindent . . . . . . . . . . . .\par
\noindent\hspace*{\parindent}Ich selbst habe die Juno nur noch zweimal beobachtet, am 26. Januar und 1. Febr[uar], am letztern Tage eigentlich nur sie gesehen, da 240 Ceti Bode, womit ich sie ein paarmal verglich, zu nördlich stand, um eine brauchbare Position daraus abzuleiten.

Die Beob. vom 26. Januar ist:
\begin{center}$1805\ \text{Jan.}\ 26.\qquad 7^{\mathrm h}1^{\mathrm m}47^{\mathrm s}\qquad 28^0\,43'\,5''\qquad 2^0\,50'\,53''.$\end{center}
\noindent Ich habe diese Beob., die sich auf 5 Vergl[eichungen] mit einem $\ast$ 8. Gr[össe] $17^0\,58'$, $3^0\,1'$ gründet, noch nicht mit den Elementen verglichen. . . . . . .
\begin{center}\small BEMERKUNGEN.\end{center}
\begingroup\small
\noindent\hspace*{\parindent}Zu I. — Der bekannte und einflussreiche Direktor der Seeberger Sternwarte v. \textsc{Zach} erhielt vom Könige von Preussen den Auftrag zu einer Gradmessung, die von der bayerischen Grenze über den Seeberg bei Gotha bis nach Westfalen gehen sollte und über die er in der Monatlichen Correspondenz, Jahrgang 1804 berichtet. Anfang August 1803 ging er auf den Brocken, von wo er fast täglich Pulversignale gab, die von anderen Punkten zum Zwecke von Längenbestimmungen beobachtet wurden. Er hatte auch \textsc{Gauss} gebeten, an diesen Beobachtungen teilzunehmen.

Zu II sind Kreismikrometerbeobachtungen von Planeten zusammengestellt, die aus Braunschweig stammen und die \textsc{Gauss} selbst nicht veröffentlicht hat.
\begin{flushright}\textsc{Brendel}.\end{flushright}
\endgroup
\begin{center}\rule{.24\linewidth}{.5pt}\end{center}
\clearpage\endgroup
